%Transcranial focused ultrasound stimulation (tFUS) is being increasingly investigated as a non-invasive approach to interrogate brain circuits and treat psychiatric disorders marked by aberrant brain activity. However, the mechanisms of tFUS, as well as its effects on large-scale brain networks, remain unclear. Here, we asked whether tFUS can reliably modulate functional connectivity (FC) with a deep brain region: the subgenual anterior cingulate cortex (sgACC), a clinically relevant area implicated in mood disorders. Using functional magnetic resonance imaging (fMRI), we recorded blood-oxygen-level dependent (BOLD) activity before, during, and after 5 minutes of tFUS targeting the sgACC. Relative to sham stimulation, we found a statistically significant increase in FC with the target both during and immediately after sonication. The increase manifested as a slow increase that was observed in approximately 80\% of the sgACC's connections, suggesting enhanced integration of the sgACC into large-scale networks. These findings support the feasibility of tFUS as a tool for modulating activity in deep brain regions with potential therapeutic applications. 

\noindent Transcranial focused ultrasound stimulation (tFUS) can modulate human brain activity, but its impact on large-scale functional connectivity (FC) \emph{in vivo} remains incompletely characterized. Here we asked whether brief tFUS to the subgenual anterior cingulate cortex (sgACC) modulates its FC with other brain regions. In a sham-controlled, within-subject study of healthy adults ($N=16$), we measured resting-state blood-oxygen-level dependent (BOLD) activity with functional magnetic resonance imaging (fMRI) before, during, and after five minutes of tFUS (i.e., five 20-second sonications followed by 40 second pauses). Subject-level sgACC–whole-brain FC was analyzed with a linear mixed-effects model (fixed: time, condition, time×condition; random: subject). tFUS increased sgACC-whole brain connectivity after sonication relative to sham. A smaller, non-significant trend was observed during sonication. Baseline-controlled analyses were employed to rule out the possibility that the observed increase in FC was due to a ``regression to the mean'' effect. Interestingly, network-level analysis revealed a clear disassociation in the evolution of resting-state FC across the 15 minute recordings: during sham sessions, the sgACC's connectivity with the default mode network (DMN) increased, while FC with other networks was largely unchanged. On the other hand, active tFUS led to a notable increase in FC with non-DMN networks (especially the cognitive control network), while connectivity with the DMN was stable. These results provide whole-brain evidence that tFUS reconfigures network dynamics in the human brain, motivating clinical applications to disorders characterized by dysregulated brain activity.

